% ===================== CONCLUSIONES =====================
\section{Conclusiones}

\begin{enumerate}
	\item Se determinó que para producir 100 t/h de vapor sobrecalentado a 106 barg y 545$^{\circ}$C, se requieren \textbf{37.67 t/h de bagazo} con 48\% de humedad y 10\% de cenizas (AR).
	\item El ratio resultante es de \textbf{2.655 toneladas de vapor por tonelada de bagazo}, valor conservador acorde con las condiciones industriales más exigentes.
	\item El alto contenido de cenizas (10\% AR) se justifica técnicamente por la incorporación de materia extraña mineral durante la cosecha mecanizada y la ausencia de lavado de caña, condiciones prevalentes en la operación de ingenios colombianos.
	\item El PCI obtenido de 6.51 MJ/kg es consistente con el rango reportado por Hugot (6.0--8.5 MJ/kg) para bagazo de alta humedad.
	\item La eficiencia del 94\% y la purga del 2\% son valores estándar de ingeniería que proporcionan un margen de seguridad adecuado para el dimensionamiento del sistema de alimentación de combustible.
\end{enumerate}
